\chapter{Results from participants}
\label{ch:participant}

\begin{center}
\fbox{\parbox{0.93\linewidth}{\centering\bfseries
No participant data has been collected.\\[4pt]
\normalfont\small
This chapter is deliberately complete in structure and empty of results. Every
table and figure that the study in \cref{ch:study} would produce is set out
below, with the analysis that fills it named in the caption, so that the
chapter can be filled in the order the data arrives rather than rebuilt.
\Cref{ch:discussion} states what has to be resolved before the sessions can
run; the first item is two soldered joints.}}
\end{center}

\section{Participants}

\begin{table}[htbp]
  \centering
  \small
  \caption[Participants]{Who took part, reported separately by role because
  the two roles are not equivalent: only the wearer carries the load. No
  identifying field is recorded.}
  \label{tab:demographics}
  \begin{tabular}{@{}lccc@{}}
    \toprule
    & \textbf{Drivers} & \textbf{Wearers} & \textbf{Pairs} \\
    \midrule
    Number & -- & -- & -- \\
    Age, median and range & -- & -- & -- \\
    Previous experience of teleoperation & -- & -- & -- \\
    Handedness & -- & -- & -- \\
    Sessions completed in full & -- & -- & -- \\
    Sessions ended early, by reason & -- & -- & -- \\
    \bottomrule
  \end{tabular}
\end{table}

\section{Task performance}

\begin{figure}[htbp]
  \centering
  \screenshot{56mm}{Completion time by task and condition. Four panels, one
  per task, each with one box per condition. The pick-and-place and two-arm
  carry panels carry fewer conditions than the reaching panel, and the missing
  conditions should be drawn as an explicit gap with a note, not omitted, so
  that a reader does not read them as data that was lost.}
  \caption[Completion time]{Completion time. To be produced by the autonomy
  level analysis. Report the omnibus test and the corrected pairwise
  comparisons, with raw and corrected values side by side.}
  \label{fig:completion}
\end{figure}

\begin{table}[htbp]
  \centering
  \small
  \caption[Task outcomes]{Success rate and completion time by task and
  condition. Cells marked \emph{n/a} are the combinations the platform cannot
  perform, explained in \cref{ch:results}.}
  \label{tab:outcomes}
  \begin{tabular}{@{}lcccc@{}}
    \toprule
    & \multicolumn{2}{c}{\textbf{Success rate}} &
      \multicolumn{2}{c}{\textbf{Time (s), median}} \\
    \cmidrule(lr){2-3}\cmidrule(lr){4-5}
    \textbf{Task} & Direct & Assisted & Direct & Assisted \\
    \midrule
    T1 reaching & -- & -- & -- & -- \\
    T2 pick and place & n/a & n/a & n/a & n/a \\
    T3 two-arm carry & n/a & n/a & n/a & n/a \\
    T4 handing over & -- & -- & -- & -- \\
    \bottomrule
  \end{tabular}
\end{table}

\begin{figure}[htbp]
  \centering
  \screenshot{50mm}{Path efficiency and clutch use. Two panels: the ratio of
  the hand's travelled distance to the straight-line distance, and the number
  of times the driver disengaged and re-engaged, both by condition.}
  \caption[How the driver worked]{Path efficiency and re-engagement count.
  These two together say whether assistance made the movement more direct or
  merely slower.}
  \label{fig:efficiency}
\end{figure}

\section{Workload and exertion}

\begin{table}[htbp]
  \centering
  \small
  \caption[Workload]{Task load index, by role and condition. Reported for both
  participants because a reduction for one is not a reduction for the other,
  and the comparison between the two columns is the point of the study.}
  \label{tab:tlx}
  \begin{tabular}{@{}lcccccc@{}}
    \toprule
    & \multicolumn{3}{c}{\textbf{Driver}} &
      \multicolumn{3}{c}{\textbf{Wearer}} \\
    \cmidrule(lr){2-4}\cmidrule(lr){5-7}
    \textbf{Subscale} & C1 & C2 & C3 & C1 & C2 & C3 \\
    \midrule
    Mental demand & -- & -- & -- & -- & -- & -- \\
    Physical demand & -- & -- & -- & -- & -- & -- \\
    Temporal demand & -- & -- & -- & -- & -- & -- \\
    Performance & -- & -- & -- & -- & -- & -- \\
    Effort & -- & -- & -- & -- & -- & -- \\
    Frustration & -- & -- & -- & -- & -- & -- \\
    \midrule
    Overall & -- & -- & -- & -- & -- & -- \\
    \bottomrule
  \end{tabular}
\end{table}

\begin{figure}[htbp]
  \centering
  \screenshot{46mm}{Exertion rating from the wearer, plotted against block
  number rather than against condition, with the condition marked on each
  point. If exertion rises monotonically with block regardless of condition,
  it is fatigue and must be treated as a covariate rather than as an effect.}
  \caption[Exertion]{Physical exertion reported by the wearer between blocks.}
  \label{fig:exertion}
\end{figure}

\section{Trust and ownership}

\begin{table}[htbp]
  \centering
  \small
  \caption[Trust and ownership]{Trust, by role and condition, and the
  wearer's sense of ownership at the end of the session. The prediction from
  the existing literature is that trust falls as the machine takes over more,
  and this study is designed to be able to reproduce or contradict that.}
  \label{tab:trust}
  \begin{tabular}{@{}lccc@{}}
    \toprule
    & \textbf{C1 direct} & \textbf{C2 assisted} & \textbf{C3 autonomous} \\
    \midrule
    Trust, driver & -- & -- & -- \\
    Trust, wearer & -- & -- & -- \\
    Ownership, wearer & -- & -- & -- \\
    Agency, wearer & -- & -- & -- \\
    \bottomrule
  \end{tabular}
\end{table}

\section{Safety events}

\begin{table}[htbp]
  \centering
  \small
  \caption[Safety during sessions]{Everything the safety layer did during the
  sessions. These are recorded automatically and should be reported whether or
  not anything happened, because a table of zeros is a result.}
  \label{tab:safety-events}
  \begin{tabular}{@{}lccc@{}}
    \toprule
    \textbf{Event} & \textbf{C1} & \textbf{C2} & \textbf{C3} \\
    \midrule
    Emergency stops, by whom & -- & -- & -- \\
    Clearance holds & -- & -- & -- \\
    Closest approach recorded, minimum & -- & -- & -- \\
    Trials invalidated by a fault & -- & -- & -- \\
    Sessions stopped at a participant's request & -- & -- & -- \\
    \bottomrule
  \end{tabular}
\end{table}

\section{What participants said}

\noindent\fbox{\parbox{0.95\linewidth}{\small\itshape
To be filled from the debrief. Free comments to be reported by role and
grouped by theme, with the number of pairs raising each. Comments from the
wearer should be reported at least as fully as comments from the driver: the
wearer is the participant with the least control over what happened and the
most direct experience of it.}}
